\section{数值语义与边界}
默认范围为非负数据，条形从零开始，长度遵循
\[
  L=W\min\!\left(1,\frac{v}{M}\right),
  \qquad v\geq 0,\quad M>0,
\]
其中 $v$ 是输入值，$M$ 是 \key{max}，$W$ 是 \key{width}。
标签的舍入与格式不改变条长。

\begin{example}{示例 3：原始数值、单位与计算占比}
\gradbar[max=200]{50}

\gradbar[max=200,
  unit={\,ms}]{50}

\gradbar[max=200,
  value format=percent]{50}
\end{example}

\key{unit} 只添加后缀，不换算数值。
\key{value format=percent} 才计算 $100v/M$，自动添加百分号，并忽略 \key{unit}。
如果原始数据已经是 86.5 这样的百分数，应使用 \key{max=100,unit={\%}}。

\begin{example}{示例 4：已有百分数与自定义标签}
\gradbar[unit={\%}]{86.5}

\gradbar[max=10,
  text={8 / 10}]{8}

\gradbar[max=10,
  text={},precision=0]{8}
\end{example}

\begin{example}{示例 5：零值和超出最大值}
\gradbar{0}

\gradbar{100}

\gradbar[max=100,
  value format=percent]{125}
\end{example}

零值只有空轨道，不会被强制画出最短长度。很小的正值也不会人为放大；
实际可见精度受 TeX 长度分辨率限制。
超出最大值时默认发出警告，并截断彩色区域，标签仍显示真实值。

\note{\key{overflow=error} 会把越界改成明确报错。
默认范围不接受负数；需要双向条时设置负的 min。非数字、零或负的最大值同样会报错。}

\clearpage

\subsection{科学计数输入与有限精度}
普通小数与 \key{1.2e-3}、\key{-2.5E4} 共用数值解析，可用于图形、CSV、筛选排序、排名、汇总和尺度。
零值有效，非零数值绝对值需位于 $10^{-100}$ 至 $10^{100}$；显式范围和显示统计值也适用这一边界。
小数点两侧各至多 101 位数字，指数部分至多三位；内部约有 16 位有效数字，更多输入位数不能保证可区分。
不接受地区千位分隔、小数逗号、公式、输入单位、NaN 或无穷大。显示精度不会改变计算和排名。
\begin{example}{科学计数输入仍保留原始数值}
\gradbar[min=0,max=1e-2,number format=scientific,precision=2]{1.25E-3}
\end{example}
